%--------------------------------------------------------------------------
\chapter{Introduction}
%--------------------------------------------------------------------------

%--------------------------------------------------------------------------
\section{Motivation}
%--------------------------------------------------------------------------

Due to their potentially low thermal conductivities, 
disordered materials (e.g., alloys, amorphous solids, aerogels) 
are used in 
applications such a photoelectric energy conversion, 
thermoelectric energy conversion, and thermally insulating barriers 
\cite{graebner_phonon_1986,cahill_lattice_1988,lu_thermal_1992,
chen_recent_2003,clarke_thermal_2005,snyder_complex_2008,
minnich_bulk_2009,toberer_phonon_2011,zebarjadi_perspectives_2012,
schiffres_gas_2012,duda_exceptionally_2013}. 
The focus of this work is on dielectric or semiconducting solids,  
where the heat is conducted by the atomic vibrational modes. 
Thermal management engineering (i.e., heat spreaders and thermal insulators) 
depends critically on accurate predictions of the vibrational mode 
properties. 
The mode properties allow 
for a deeper understanding of the nature of thermal transport at the 
atomistic scale, which can be used to understand the 
effects of vibrational scattering mechanisms in bulk and 
nanostructured materials, which is an area of 
active research \cite{chen_recent_2003,snyder_complex_2008,
turney_predicting_2009,minnich_bulk_2009,toberer_phonon_2011,
mcgaughey_size-dependent_2011,zebarjadi_perspectives_2012,
mcgaughey_nanostructure_2012,regner_broadband_2013}. 

Thermal transport in ordered (crystalline) materials is well understood 
in terms of the phonon gas model 
\cite{maradudin_dynamical_1974,wallace_thermodynamics_1972,
srivastava_physics_1990,dove_introduction_1993,ziman_electrons_2001}. 
Predicting the thermal conductivity requires the properties of the 
full spectrum of vibrational modes. 
Theoretical formulations for phonon properties date back to the work of 
Callaway \cite{callaway_model_1959}, 
Holland \cite{holland_analysis_1963}, 
Klemens \cite{klemens_scattering_1955}, 
and Slack \cite{slack_thermal_1979}. Their theory derived and/or 
postulated analytical models for the phonon dispersion and lifetimes 
(typically based on low-frequency limits) to be used in a solution 
of the BTE for predicting thermal conductivity. By fitting the
BTE solution to experimental bulk thermal conductivity data, expressions 
for the lifetimes are obtained. While good fits to the experimental 
data can be found, this agreement may be due to the numerous fitting 
parameters and not due to the correctness of the phonon lifetime models 
\cite{chung_role_2004}. 
Atomistic techniques, which can predict the properties of individual phonon
modes without any assumptions about the dispersion or scattering, 
are thus required. Newly-developed computational methods based on density 
functional theory (DFT) calcualtions can make experimentally-accurate 
predcitions of the phonon properties and thermal conductivity of bulk 
crystals, whose structures are based on a unit cell with a small number of 
atoms \cite{broido_intrinsic_2007,ward_ab_2009,ward_intrinsic_2010,
esfarjani_heat_2011,lindsay_thermal_2012,shiga_microscopic_2012,
luckyanova_coherent_2012,lindsay_first-principles_2013}. 

Disordering a crystal (i.e., alloying, 
amorphization) 
breaks down the phonon gas theory and a new theoretical formulation 
is necessary. Theoretical formulations for disordered crystals begin  
with perturbation theory 
\cite{klemens_scattering_1955,abeles_lattice_1963,
wallace_thermodynamics_1972,maradudin_dynamical_1974,
tamura_isotope_1983}, 
which is only valid for weakly disordered 
systems (i.e., isotopically disordered crystals) and bases the 
calculations on a unit cell with a small 
number of atoms. 
The study of disordered lattices (i.e., alloys) has become 
increasingly quantitative as researchers seek lower thermal 
conductivity thermoelectric materials 
\cite{chen_recent_2003,snyder_complex_2008,toberer_phonon_2011,
shiomi_thermal_2011,
zebarjadi_perspectives_2012,tian_phonon_2012,li_thermal_2012}. 
Modeling can be experimentally accurate using 
DFT calcualtions together with perturbative 
methods 
\cite{broido_intrinsic_2007,ward_ab_2009,ward_intrinsic_2010,
garg_role_2011,
li_thermal_2012,lindsay_thermal_2012,lindsay_first-principles_2013}. 
Because DFT calculations are computationally expensive, they rely on 
the perturbative methods. 
The comparison of the perturbative DFT predictions to experiments 
is only qualititative for some disordered materials, particularly 
for those with lower thermal conductivities 
\cite{kundu_role_2011,shiomi_thermal_2011,tian_phonon_2012,
tian_enhancing_2012}. 
Thus, there is a need to asses the 
applicability of the perturbative methods used in DFT calculations 
for disordered crystals with a range of disorder 
\cite{luckyanova_coherent_2012,kundu_role_2011,shiomi_thermal_2011,
tian_phonon_2012,tian_enhancing_2012,luckyanova_coherent_2012}.

Amorphous materials have been well-studied in the literature for 
their interesting properties such as the Boson peak 
\cite{schirmacher_harmonic_1998,schirmacher_acoustic_2007,
chumakov_equivalence_2011}, 
excess modes \cite{sette_dynamics_1998,zink_excess_2006}, 
and, most relevant to the current study, phonon-like 
thermal transport 
\cite{allen_thermal_1993,feldman_thermal_1993,freeman_thermal_1986,
graebner_phonon_1986,love_estimate_1990,cahill_thermal_1994,
feldman_numerical_1999,baldi_thermal_2008,liu_high_2009,
yang_anomalously_2010,zink_thermal_2006,
hondongwa_ultrasonic_2011,masciovecchio_evidence_2006,
baldi_sound_2010,baldi_elastic_2011,baldi_emergence_2013,
he_heat_2011,pompe_thermal_1988,
hasselman_thermal_1989,kuo_thermal_1992,
wada_thermal_1996,moon_thermal_2002,zink_excess_2006,
zeller_thermal_1971,cahill_lattice_1988,cahill_heat_1989,
lee_heat_1997,yamane_measurement_2002,
regner_broadband_2013}.  
Theoretical modeling of amorphous materials has previously relied 
on phenomenological models, which have limited predictive capability 
\cite{kittel_interpretation_1949,slack_thermal_1979,freeman_thermal_1986,
graebner_phonon_1986,cahill_heat_1989,cahill_lattice_1988}. 
The theoretical formulation 
for the mode properties in amorphous materials has only been 
recently developed.
\cite{allen_thermal_1993,feldman_thermal_1993,allen_diffusons_1999} 
Like a disordered crystal, an amorphous material has a unit cell 
with an infinite number of atoms, by definition. 
However, modeling of amorphous materials is  
difficult because the topological 
disorder prohibits the use of perturbative methods. Instead, the 
disordered material must be modeled explicitly. 
A finite unit cell can be used to model an amorphous material, but 
the required size of these models is typically large 
\cite{feldman_thermal_1993,feldman_numerical_1999}, which 
limits the use of DFT calculations. As with disordered crystals, 
modeling of amorphous materials is becoming increasingly 
quantitative as new experimental techniques for measuring 
thermal properties reach higher levels of resolution 
\cite{masciovecchio_evidence_2006,baldi_thermal_2008,
baldi_sound_2010,baldi_emergence_2013,baldi_elastic_2011,
liu_high_2009,yang_anomalously_2010,
regner_broadband_2013,sultan_heat_2013}. 
Thus, there is a need to asses the predictive capabilities of 
the theoretical models for predicting vibrational mode properties 
in amorphous materials.

%--------------------------------------------------------------------------
\section{\label{ThermalModes}
Vibrational Modes: Phonons, Propagons, Diffusons, and Locons}
%--------------------------------------------------------------------------

For a perfect lattice, all vibrational modes are phonon modes, which by 
definition are delocalized, propagating plane waves 
\cite{ziman_electrons_2001}. For disordered materials, all 
vibrational modes are not phonons because there is no translational 
order.  Allen and Feldman define all disordered vibrational modes to be 
vibrons because they are all bosons and follow Bose-Einstein 
statistics \cite{allen_diffusons_1999}. The vibrons are propagons, 
diffusons, or locons. Propagons are propagating and delocalized 
(i.e., phonon-like) and are typically found in the low-frequency 
range of the vibrational spectrum \cite{allen_diffusons_1999}. 
Diffusons are non-propagating and delocalized modes that couple 
harmonically due to the disorder and spatial delocalization 
\cite{allen_thermal_1993}. While diffusons are non-propagating, 
they can contribute a significant amount to thermal transport 
\cite{feldman_thermal_1993,feldman_numerical_1999}. 
Locons are non-propagating and spatially 
localized modes that not contribute to thermal transport, 
and are typically found at the high frequency of the vibrational 
spectrum \cite{allen_thermal_1993,allen_diffusons_1999}.

The spectrum of vibrons is depcited in Fig. 
\ref{F:allen_diffuson_1999_fig1} for a model of amorphous silicon
\cite{allen_diffusons_1999}. The propagons exist at low frequencies below 
the Ioffe-Regel (IR) limit \cite{taraskin_determination_1999}, 
which is the proposed transition between 
propagating and non-propagating modes. The IR limit is discussed in 
Sections \ref{S:Phonon Lifetimes} and \ref{S:Life}. 
Propagons and diffusons are called extendons because they are 
delocalized. The mobility edge marks the transition from 
the delocalized diffusons to the localized locons 
\cite{allen_diffusons_1999}.
 
%--------------------------------------------------------------------------
\begin{figure}
\begin{center}
\centering
\includegraphics[scale=0.4]
{/home/jason/thesis/thesis/allen_diffuson_1999_fig1.eps}
\end{center}
\caption{\label{F:allen_diffuson_1999_fig1} 
Density of vibrational states for a model of amorphous silicon and 
classification of vibrational modes (vibrons). 
Figure is reproduced from Ref. \citenum{allen_diffusons_1999}.
}
\end{figure}
%--------------------------------------------------------------------------
\clearpage

%--------------------------------------------------------------------------
\section{\label{S-SystemMode}
System-level versus Mode-level Predictions}
%--------------------------------------------------------------------------

%--------------------------------------------------------------------------
\subsection{\label{S-ktopdown} 
System-level Prediction}
%--------------------------------------------------------------------------

% There are two commonly-used top-down techniques for predicting thermal 
% conductivity from MD simulations: the Green-Kubo method 
% \cite{mcquarrie2000,mcgaughey2006c}, which uses equilibrium simulations, 
% and the direct method \cite{schelling2002,mcgaughey2006c}, which uses 
% steady-state non-equilibrium simulations. These two methods are compared 
% in Table \ref{T-ktopdown}. While so-called ``quantum corrections" to 
% thermal conductivities predicted from classical MD simulations have been 
% proposed, Turney et al. demonstrate that these corrections are not 
% rigorous and should not be applied \cite{turney2009b}. Comparison of 
% MD-predicted thermal conductivities to experimental measurements should 
% therefore be limited to high temperatures, around and above a material's 
% Debye temperature. Top-down techniques are important in the context of 
% this chapter because they can be used to validate the bottom-up 
% techniques that are our focus.
% 
% The Green-Kubo and direct methods can be applied to crystals, alloys, 
% amorphous solids, and fluids. Recent work has focused on computational 
% aspects, notably accounting for size effects, quantifying and reducing 
% uncertainty, and reducing computational time 
% \cite{landry2008,zhou2009,jchen2010,
% sellan2010a,howell2011a,howell2011b,
% he2012,howell2012,jones2012}.
% The Green-Kubo method relates the equilibrium fluctuations
% of the heat current vector, $\mathbf{S}$, to the thermal conductivity 
% tensor
% via the fluctuation-dissipation theorem. The thermal conductivity in 
% direction $\alpha$ is given by \cite{mcquarrie2000}
% \begin{eqnarray}\label{E-GK}
% k_{\alpha} = \f{1}{k_\mathrm{B} V T^2} 
% \int_0^{\infty} \lr{S_{\alpha}(t)S_{\alpha}(0)}dt, 
% \end{eqnarray}
% where $k_{\mathrm B}$ is the Boltzmann constant, $t$ is time, and 
% $S_{\alpha}$ and $\lr{S_{\alpha}(t)S_{\alpha}(0)}$ are the $\alpha$ 
% components of the heat current vector and the heat current autocorrelation 
% function (HCACF). The heat current vector is calculated using the atomic 
% positions and velocities. When using the Green-Kubo method, we recommend 
% that the MD simulations be run in the $NVE$ ensemble.\footnote{As mentioned 
% in Section \ref{S-MDformulation}, $NVE$ simulations are advantageous 
% because the equations of motion are not modified. The Green-Kubo method 
% formula for thermal conductivity, Eq.~\eqref{E-GK}, however, is derived 
% in the $NVT$ ensemble, leading to the need to specify the temperature of 
% the MD simulation. Temperature is not well-defined in the $NVE$ ensemble, 
% but can be approximated using the system kinetic energy and an assumption 
% of equipartition of energy (very good for classical MD simulation of 
% solids).} The Green-Kubo method is advantageous for anisotropic materials 
% because it can be used to predict the full thermal conductivity tensor.

There are several atomistic methods for predicting the thermal 
conductivity at the system level. The molecular dynamics (MD)-based 
\href{http://lammps.sandia.gov/doc/fix_thermal_conductivity.html}
{direct method} predicts the thermal conductivity by a direct 
application of the 
\href{http://en.wikipedia.org/wiki/Thermal_conduction#Fourier.27s_law}
{Fourier law} for heat conduction 
\cite{skye_thermal_2008,landry_complex_2008,
landry_effect_2009,landry_thermal_2009,
turney_predicting_2009,thomas_thermal_2010}. 
While the direct method 
is conceptually simple, it can be difficult to account for finite-size 
effects. The 
\href{http://lammps.sandia.gov/doc/compute_heat_flux.html}{Green-Kubo}
method can predict the thermal conductivity from equilibrium MD 
simulations. These equilibrium simulations are often simpler 
to setup and account for finite-size effects 
(see Section \ref{Appendix_A:Finite}).  
The Green-Kubo and direct methods can be applied to crystals, alloys, 
amorphous solids 
\cite{mcgaughey_quantitative_2004,mcgaughey_thermal_2004,
landry_complex_2008,he_lattice_2012,
he_heat_2011,he_morphology_2011,he_thermal_2011,turney_predicting_2009,
landry_complex_2008,shiomi_thermal_2011,esfarjani_heat_2011}, 
and is described in detail elsewhere 
\cite{mcquarrie_statistical_2000,mcgaughey_phonon_2004}.
While the GK and direct methods can take into 
account the effects of disorder, 
they predict no mode-level 
properties of the thermal carriers. 

While so-called ``quantum corrections" to 
thermal conductivities predicted from classical MD simulations have been 
proposed, Turney et al. demonstrate that these corrections are not 
rigorous and should not be applied \cite{turney_assessing_2009}. 
Comparison of MD-predicted thermal conductivities to experimental 
measurements should 
therefore be limited to high temperatures, around and above a material's 
Debye temperature. The GK method, which is a system-level technique, 
is important for the present study because it can be used to validate 
the mode-level techniques that are the focus. The capabilities 
of the GK method are summarized in Table \ref{T-comparison}. It is 
important to note that the theoretical formulation of the GK method 
prohibits its use with DFT calculations, where the per atom potential 
energies cannot be rigorously defined.\cite{martin_electronic_2004}

%--------------------------------------------------------------------------
\subsection{\label{S-kbottomup}Mode-level Prediction}
%--------------------------------------------------------------------------

While the total thermal conductivity is typically the quantity of 
interest for modeling and engineering analysis, the mode-level 
properties can be just as important. The mode-level properties 
allow for a deeper understanding of the total thermal conductivity.  
The mode-level properties are used to understand the difference 
between predictions for alloys in Chapter \ref{chapter:vc} 
and amorphous materials in Chapter \ref{chapter:mfp}. 

The Callaway-Holland (CH) theory, which is based on the work of 
Callaway \cite{callaway_model_1959}, Holland \cite{holland_analysis_1963},
Klemens \cite{klemens_scattering_1955}, 
and Slack \cite{slack_thermal_1979}, uses analytical models for phonon 
dispersion and lifetimes. These analytical models are often empirically 
fit to experimental results, which limits the predictive capability of 
the CH theory. While the CH theory can calculate mode properties 
from the analytical models, it is unclear how accurate these properties 
are.\cite{chung_role_2004}

There are several methods that can predict the mode-level vibrational 
properties. 
Anharmonic lattice dynamics (ALD) calculations 
\cite{turney_predicting_2009,esfarjani_heat_2011}, 
based on perturbation theory 
\cite{maradudin_dynamical_1974,srivastava_physics_1990},
can be used along with phonon-defect scattering perturbation 
theory to predict the mode properties and thermal conductivity 
of mass and/or bond disordered lattices (see Section \ref{S:From VC-ALD}) 
\cite{klemens_scattering_1955,klemens_thermal_1957,mattis_phonon_1957,
tamura_isotope_1983}.  
However, ALD is a perturbative method that is valid at low 
temperatures where 3-phonon processes dominate 
\cite{wallace_thermodynamics_1972,maradudin_dynamical_1974,
turney_predicting_2009}. 
Quantum or classical statistics can be used with the ALD method. 
Because the ALD method is based on a unit cell 
and disorder is included as a perturbation, the mode properties 
predicted are for phonons. One of the objectives of this work is 
to investigate if the phonon properties predicted by the ALD method 
are consistent with the properties of propagons in a fully 
disordered system (see Chapter \ref{chapter:vc}). 

The inputs to the ALD method are the harmonic and anharmonic 
force constants from the interatomic interactions, 
which can be obtained from empirical interatomic potentials or 
DFT calculations \cite{esfarjani_method_2008,turney_predicting_2009}. 
Because the ALD method is computationally inexpensive, it can be 
used with computationally expensive DFT calculations to predict 
thermal conductivities that are experimentally accurate 
\cite{broido_intrinsic_2007,ward_intrinsic_2010,ward_ab_2009,
garg_role_2011,lindsay_thermal_2012,li_thermal_2012,lindsay_thermal_2012,
lindsay_first-principles_2013,li_thermal_2012}. 
However, most studies using ALD and DFT have been limited to 
disordered lattices (i.e., alloys) whose thermal conductivities 
are large and dominated by low-frequency propagons 
\cite{broido_intrinsic_2007,ward_intrinsic_2010,ward_ab_2009,
garg_role_2011,lindsay_thermal_2012,lindsay_thermal_2012,
lindsay_first-principles_2013}. 
Because no comprehensive study has been performed on a range 
of materials, it is unclear what limitations exist for using the ALD and 
phonon-defect perturbative method 
\cite{luckyanova_coherent_2012,kundu_role_2011,shiomi_thermal_2011,
tian_phonon_2012,tian_enhancing_2012,luckyanova_coherent_2012}. 
The limitations of the ALD and phonon-defect perturbative 
method is investigated in Chapter \ref{chapter:vc}. 

The key to explicitly incorporating the effects of disorder 
is to use large disordered supercells, or a unit cell with a large 
number of atoms. The use of computationally 
inexpensive empirical potentials allows for calculations to be 
performed on a large number of atoms, up to $10^5-10^6$ 
\cite{he_heat_2011,he_thermal_2011,he_morphology_2011,
he_lattice_2012}. 
The Allen-Feldman (AF) theory computes 
the contribution of diffusons 
to thermal conductivity from large disordered supercells 
\cite{allen_thermal_1993}. 
The AF theory calculations are based on harmonic lattice dynamics that 
rely on the system disorder to couple modes thermally 
\cite{allen_thermal_1993,feldman_thermal_1993,feldman_numerical_1999}.  
However, because the AF theory is harmonic, it is not suited to studying 
the low-frequency propagons in disordered materials 
\cite{feldman_thermal_1993}, particularly for disordered 
lattices where the harmonic theory predicts Rayleigh scattering, which 
causes a divergent thermal conductivity with increasing system 
size \cite{sheng_heat_1991,sheng_introduction_2006,vitelli_heat_2010}. 
The AF theory can use quantum or classical expressions 
for the specific heat. While the AF theory should be limited to low 
temperatures where the harmonic approximation is valid, 
anharmonic effects have been shown to be 
small for disordered solids such as amorphous silicon 
\cite{feldman_thermal_1993}.

The normal mode decomposition (NMD) method can be used to predict the 
lifetime of all vibrons from MD simulations 
\cite{bickham_calculation_1998,
bickham_numerical_1999,fabian_numerical_2003,
ladd_lattice_1986,mcgaughey_quantitative_2004,donadio_thermal_2007,
henry_spectral_2008,donadio_atomistic_2009,turney_predicting_2009,
goicochea_thermal_2010,he_morphology_2011,he_thermal_2011,he_heat_2011,
he_lattice_2012,sosso_thermal_2012,qiu_molecular_2012,wang_mode-wise_2013}.  
While this method can only predict the 
mode lifetimes, these can be used with other predictive methods 
and assumptions to predict the thermal conductivity (see Sections 
\ref{S:Diffusivities_vc} and \ref{S:Diffusivities_mfp}). 
Because the lifetimes predicted by NMD are obtained from MD simulations, 
they should be used in a classical limit \cite{turney_assessing_2009}. 
The thermal conductivity predicted 
by NMD begins to deviate from GK and direct-method prediction 
near about half the melting temperature for Lennard-Jones argon 
\cite{turney_predicting_2009}. 

Because of the theoretical limitations of the methods discussed above, 
no one method is sufficient for predicting the mode 
properties of all vibrons in a given disordered material. The 
theoretical methods are summarized and compared in 
Table \ref{T-comparison}. 
The methods in Table \ref{T-comparison} are ranked in the order 
of their thermal conductivity predictive capability based on 
prior work in the literature, starting 
with the least predictive on the left. 
For example, the CH theory relies completely on empirical fitting, 
while the GK method can predict the thermal conductivity of some 
materials in any phase (crystalline, amorphous, liquid) 
\cite{mcgaughey_phonon_2004}. 
In Chapter \ref{Conclusion}, the predictive capabilities are re-ranked in 
Table \ref{T-comparison-predictive} based on their ability to predict the 
mode properties in disordered systems from the results in this work.

%--------------------------------------------------------------------------
\begin{center}
%\squeezetable
\begin{table}
\small
\caption{\label{T-comparison}Comparison of theoretical 
techniques for predicting vibrational mode properties and thermal 
conductivity for disordered systems.}
%\begin{tabular}{llllll}
\begin{tabular}{p{0.9in}|p{0.9in}|p{0.9in}|p{0.9in}|p{0.9in}|p{0.9in}}
\hline\hline
&CH Theory 
&ALD
&AF Theory
&NMD
&GK \\ 
\hline\hline
Description
&Empirical
&Predictive 
&Predictive 
&Predictive
&Predictive\\
\hline
Input
&Fitted Models
&Force constants \newline (empirical potential or DFT)  
&Force constants \newline (empirical potential or DFT)
&Forces and Force Constants \newline (empirical potential or DFT) 
&Forces \newline (empirical potential) \\
\hline
Disorder
&As additional \newline perturbation 
&As additional \newline perturbation 
&Naturally included 
&Naturally included 
&Naturally included \\
\hline
Statistics
&Bose-Einstein or \newline Boltzmann 
&Bose-Einstein or \newline Boltzmann 
&Bose-Einstein or \newline Boltzmann
&Boltzmann 
&Boltzmann\\
\hline
Temperature
&All 
&Low 
&All
&Mid-range
&All\\
\hline
Modes
&Phonons (Vibrons?)
&Phonons (Vibrons?)
&Diffusons
&Phonons and Vibrons
&N/A\\
\hline \hline
\end{tabular}
\end{table}
\end{center}
%--------------------------------------------------------------------------
\clearpage

%--------------------------------------------------------------------------
\subsection{\label{S-CompCost}
Computational Cost vesus Predictive Ability}
%--------------------------------------------------------------------------

The computational costs of the methods discussed in Sections 
\ref{S-ktopdown} and \ref{S-kbottomup}  
are given in Table \ref{T-comparison-cost}.  
In most cases, the 
computational cost of a method is known $\emph{a priori}$ and should 
be considered before any analysis is performed. The computational 
costs for these methods can be represented by the total number of 
atoms, $N_a$, that are being simulated. For ordered and disordered 
systems, the total 
number of atoms can be related to the number of atoms in the unit 
cell, $n$, and the 
number of unit cells, $N$ ($N$ = 1 for a disordered system), 
which gives $N_a = N n$ (see Section \ref{Subsection_Comp_Details_1}).

One of the goals 
of this work is to thoroughly asses the predictive 
abilities of these methods in relation to their computational costs. 
For example, 
The NMD method can predict the mode lifetimes for the vibrational 
modes in ordered and disordered systems (i.e., phonons and vibrons), 
making its predictive ability high 
\cite{ladd_lattice_1986,mcgaughey_quantitative_2004,donadio_thermal_2007,
henry_spectral_2008,donadio_atomistic_2009,turney_predicting_2009,
goicochea_thermal_2010,he_morphology_2011,he_thermal_2011,he_heat_2011,
he_lattice_2012,sosso_thermal_2012,qiu_molecular_2012,wang_mode-wise_2013}. 
However, the explicit use of disorder makes the NMD method's 
computational cost high (see Section \ref{Subsection_Comp_Details_4}), 
making it particularly challenging to use with DFT calculations 
\cite{koker_thermal_2009,huang_filler-reduced_2010,
huang_ab_2008,luo_molecular_2011,esfarjani_heat_2011,
shiomi_thermal_2011}. 

For the AF Theory and NMD method, the major computational cost is the 
eigenvalue solution of the dynamical matrix 
(see Sections \ref{S:VC Gamma DOS} and \ref{S:Simulation}), which scales 
poorly in both CPU and memory costs. 
For ALD, $N_a$ determines the resolution of the 
\href{https://en.wikipedia.org/wiki/Brillouin_zone}
{Brillouin zone} sampling (see Section \ref{Subsection_Comp_Details_1} 
and Ref. \citenum{turney_predicting_2009-1}). Some systems (such as 
the silicon alloys studied in Section \ref{S:SW}) 
require a fine sampling 
of the BZ, which increases the ALD computational costs. Using large 
systems can be avoided by following the extrapolation procedure 
which is discussed in Appendix \ref{Appendix_A:Finite}.
Efficient MD 
codes like LAMMPS scale linearly with the number of atoms in 
the system which makes the GK method (see Section 
\ref{S:Thermal Conductivity}) 
computationally-inexpensive.
 
The computational costs for each method are ranked in 
Table \ref{T-comparison-cost} starting with least 
expensive on the left. 
The rankings do not take into account the 
coefficients of the computational 
scalings. For example, to perform NMD on a system 
of 6912 atoms in less 
than 24 hours may take 100 central processing 
units (CPUs).  
%(see Appendix \ref{A:unitcell}). 
The same number of 
CPUs are required 
to perform the GK method on a system of 
100000 atoms in the same amount of 
time. Coefficients of the computational scalings would take into 
account the difference in computational cost of classical pair 
potentials versus multi-body potentials, for example. 
Thus, the ranking in Table \ref{T-comparison-cost} is to be used as a 
rough guideline. The maximum system 
sizes are those achieved in this work.  

The GK method, which scales most efficiently in both memory and CPU 
cost, has the most potential to simulate systems with increased 
$N_a$, but does not predict the properties of any modes. 
While the properties for a large number of modes 
can be predicted the ALD method, 
no explicit atomic dynamics are simulated and the calculation is based 
on a small number of atoms. The poor scaling of the 
computational cost with system size for the ALD method prohibits 
increasing the system size significantly. The same is true for the 
AF and NMD methods. Suggestions are given in 
Section \ref{Future:Timescales} for increasing the system sizes for the 
AF and NMD methods, as well as suggestion for other MD-based methods.

%--------------------------------------------------------------------------
\begin{center}
%\squeezetable
\begin{table}
\small
\caption{\label{T-comparison-cost}Ranking of computational costs 
for theoretical 
techniques for predicting vibrational mode properties and thermal 
conductivity, starting with the least expensive on the left.}
%\begin{tabular}{llllll}
\begin{tabular}{p{0.9in}|p{0.9in}|p{0.9in}|p{0.9in}|p{0.9in}|p{0.9in}}
\hline\hline
&CH Theory 
&AF Theory
&NMD
&GK
&ALD\\ 
\hline
Computational Cost\\
\hline
CPU\newline
Memory 
&Minimal \newline Minimal
&High \newline High 
&High \newline High 
&High \newline Low 
&High \newline Low \\
\hline
Maximum $N_a$
&N/A
&6912
&6912
&800000
&8\\
\hline
Maximum number of Modes
&N/A
&20736 
&20736 
&N/A
&592704\\
\hline
Computational Scaling\\
\hline
CPU\newline
Memory 
&Minimal \newline Minimal
&$(Nn)^{3}$ \newline $(Nn)^{2}$ 
&$(Nn)^{2}$ \newline $(Nn)^{2}$ 
&$(Nn)^{1}$ \newline $(Nn)^{1}$ 
&$N^{2}n^{4}$ \newline $Nn^{2}$  \\
\hline\hline
\end{tabular}
\end{table}
\end{center}
%--------------------------------------------------------------------------
\clearpage

%--------------------------------------------------------------------------
\section{\label{S-CompCost}
Overview and Scope}
%--------------------------------------------------------------------------

The objective of this work is to develop a theoretical and computational 
framework for predicting the vibrational mode properties in a range 
of disordered materials using atomistic techniques. 
This objectives will be achieved through the use of four predictive 
methods: MD-based GK, MD-based NMD, ALD, AF Theory 
(see Section \ref{S-SystemMode}). 
The analysis is performed on atomistic models of a wide-range of 
ordered and disordered systems. Emphasis is placed on investigating each 
material with as many of the four methods as possible to 
analyze the thermal properties self-consistently. 
The outline is as follows:

In Chapter \ref{Chapter:SED}, two MD-based methods for predicting 
vibrational mode properties are compared and evaluated. The vibrational 
mode properties and thermal conductivities are predicted for 
models of three materials: argon, silicon, and 
a carbon nanotunbe. Theoretical derivations are presented 
and the computational framework for performing the analysis is discussed. 
Recommendations are made based on the comparison of the two methods 
for future work using MD-based methods 
(see Section \ref{Future:Timescales:DSF}).

In Chapter \ref{chapter:vc}, the virtual crystal approximation for 
predicting the vibrational properties of alloys is investigated. 
Two model systems, Lennard-Jones argon and Stillinger-Weber
silicon, are used to perform the analysis. 
The analysis is performed 
using molecular dynamics simulations and lattice
dynamics calculations. 
The vibrational mode properties in crystalline, 
alloyed, and amorphous materials are studied 
using fully atomistic models. 
Mode frequencies and lifetimes are first 
calculated by treating the disorder explicitly and under the 
virtual crystal approximation. The two model systems 
have different vibrational mode spectra, which plays an important role in 
determining the applicability of the virtual crystal approximation. 
The mode properties are then used to predict thermal conductivity, and 
common phenomenological limits are used to interpret the results. Important 
connections between disordered lattices and amorphous materials are made 
through a comparison of their mode properties, which solidifies 
the modeling framework for studying the propagating and non-propagating 
vibrational modes in disordered systems. 

In Chapter \ref{chapter:mfp}, the properties of the propagating and 
non-propagating vibrational modes in amorphous silica and silicon 
are predicted using realistic atomistic models.   
The calculations are performed using molecular dynamics simulations, 
lattice dynamics calculations. 
Using the mode properties, 
the thermal conductivity accumulation functions are predicted and compared 
with recent experimental measurements by 
Regner et al.\cite{regner_broadband_2013} and thin film 
thermal conductivity measurements.  The two amorphous materials studied 
have significantly different contributions from propagating and 
non-propagating vibrational modes, which is evident in their 
accumulation functions and the comparison of the predictions 
to experimental measurements.

In Chapter \ref{Conclusion}, the major contributions of the work are
presented and suggestions for future study are discussed. 
The results presented in this work provide a theoretical and 
computational framework for the study of emerging disordered 
and nanoscaled systems. 

